\section{Circuit Complexity --- Circuit Sat Reduction Machine Tapes}
%%% generated by the blueprint pipeline from TCSlib.Complexity.CircuitComplexity.CircuitSatReductionMachineTapes
\subsection*{Circuit Complexity --- Circuit Sat Reduction Machine Tapes: Introduction}

Declarations of the Lean module \texttt{TCSlib.Complexity.CircuitComplexity.CircuitSatReductionMachineTapes}.

\subsection*{Circuit Complexity --- Circuit Sat Reduction Machine Tapes: Results}

\begin{definition}[Emitter machine states]
\lean{BoolCircuit.CktSatReduction.Emitter.St}
\label{BoolCircuit.CktSatReduction.Emitter.St}
\leanok
\uses{BoolCircuit.CktSatReduction.Emitter.Ph, BoolCircuit.GateKind}
The emitter machine has states for reading a phase; emitting a gate kind and one of
three tape indices at a program counter below $64$; printing a word or a bit with the
same parameters and a tape index; clearing a tape at one of the three indices; clearing
the second scratch tape; and incrementing.
\end{definition}

\begin{definition}[Tape-specific machine action]
\lean{BoolCircuit.CktSatReduction.Emitter.act}
\label{BoolCircuit.CktSatReduction.Emitter.act}
\leanok
\uses{BoolCircuit.CktSatReduction.Emitter.St, Turing.Action}
For one of the three work tapes $t$, an input-head movement, an optional write
instruction, an output symbol, and a successor state, define an action that applies the
supplied write instruction and movement to tape $t$, leaves the other two work tapes
unwritten and stationary, and records the specified input-head movement, output symbol,
and successor state.
\end{definition}

\begin{definition}[Control machine action]
\lean{BoolCircuit.CktSatReduction.Emitter.ctl}
\label{BoolCircuit.CktSatReduction.Emitter.ctl}
\leanok
\uses{BoolCircuit.CktSatReduction.Emitter.St, Turing.Action}
Given an input-head movement, an optional output Boolean symbol, and an optional
successor emitter state, define an action that leaves all three work tapes unwritten and
stationary and records the specified input-head movement, output symbol, and successor
state.
\end{definition}

\begin{definition}[Clause-cleanup next state]
\lean{BoolCircuit.CktSatReduction.Emitter.clNext}
\label{BoolCircuit.CktSatReduction.Emitter.clNext}
\leanok
\uses{BoolCircuit.CktSatReduction.Emitter.St}
For a tape index $t\in\{0,1,2\}$, define the next cleanup state to be the second cleanup
state when $t=1$, and the increment state otherwise.
\end{definition}

\begin{definition}[Emitter transition function]
\lean{BoolCircuit.CktSatReduction.Emitter.tr}
\label{BoolCircuit.CktSatReduction.Emitter.tr}
\leanok
\uses{BoolCircuit.CktSatReduction.Emitter.St, BoolCircuit.CktSatReduction.Emitter.act, BoolCircuit.CktSatReduction.Emitter.clNext, BoolCircuit.CktSatReduction.Emitter.ctl, BoolCircuit.CktSatReduction.Emitter.lstep, BoolCircuit.CktSatReduction.Emitter.prog, Turing.Action}
Define the machine’s one-step transition as follows. In a reading state, it uses the
current phase and input symbol to decide whether to halt, advance the phase, mark a work
tape, or begin emitting the gate’s Tseitin clauses. In an emission state, it processes
the next instruction in the gate’s clause program: output instructions emit their
Boolean value, while print instructions move the indicated work-tape head through the
write and boundary-check states; when the program ends, the machine clears the
designated work-tape cells and advances to the gate-reading phase. The cleanup and
increment states perform those respective tape actions and return to reading the gate
phase.
\end{definition}

\begin{definition}[Finite clause-emitting machine]
\lean{BoolCircuit.CktSatReduction.Emitter.emTM}
\label{BoolCircuit.CktSatReduction.Emitter.emTM}
\leanok
\uses{BoolCircuit.CktSatReduction.Emitter.St, BoolCircuit.CktSatReduction.Emitter.tr, Turing.FinTM}
Define a finite Turing machine over the Boolean alphabet with three work tapes, whose
states are the emitter states and whose initial state is the reading state for the
initial phase. Its transition function is the emitter transition defined by the phase
and tape behavior.
\end{definition}

\begin{theorem}[One step from an active state]
\lean{BoolCircuit.CktSatReduction.Emitter.step_some}
\label{BoolCircuit.CktSatReduction.Emitter.step_some}
\leanok
\uses{BoolCircuit.CktSatReduction.Emitter.St, BoolCircuit.CktSatReduction.Emitter.emTM, BoolCircuit.CktSatReduction.Emitter.tr, Turing.Action.apply, Turing.Cfg, Turing.Cfg.inputSymbol, Turing.MultiTapeTM.step}
\difficulty{1}
Let $w$ be a finite input word, and consider a configuration of a three-work-tape Turing
machine whose state is an active state $q$, whose input-head position is $p$, whose work
tapes are $T_0,T_1,T_2$, whose work-head positions are $h_0,h_1,h_2$, and whose output
word is $o$. The machine’s next configuration is the result of applying to this
configuration the action prescribed by its transition function at state $q$, reading the
symbol under the input head and the symbols under the three work-tape heads.
\end{theorem}

\begin{theorem}[Effect of a selected-tape action]
\lean{BoolCircuit.CktSatReduction.Emitter.apply_act}
\label{BoolCircuit.CktSatReduction.Emitter.apply_act}
\leanok
\uses{BoolCircuit.CktSatReduction.Emitter.St, BoolCircuit.CktSatReduction.Emitter.act, Turing.Action.apply, Turing.Cfg, Turing.UnaryTape.wrT, Turing.moveInputPos}
\difficulty{4}
Let a three-work-tape Turing-machine configuration have input head position $p$, work
tapes $T_0,T_1,T_2$, work-head positions $h_0,h_1,h_2$, and output string $O$. For a
selected tape $t\in\{0,1,2\}$, an action that moves the input head by $im$, prescribes
optional write $wr$ and head movement $m$ on tape $t$, emits the optional bit $o$, and
sets the state to $q$ produces the configuration with state $q$, input head at the
clamped position obtained by moving $p$ by $im$, tape $t$ updated by applying $wr$ at
$h_t$, head $t$ moved to $h_t+m$, all other work tapes and heads unchanged, and output
$O$ followed by $o$ if present.
\end{theorem}

\begin{theorem}[Effect of a control action]
\lean{BoolCircuit.CktSatReduction.Emitter.apply_ctl}
\label{BoolCircuit.CktSatReduction.Emitter.apply_ctl}
\leanok
\uses{BoolCircuit.CktSatReduction.Emitter.St, BoolCircuit.CktSatReduction.Emitter.ctl, Turing.Action.apply, Turing.Cfg, Turing.moveInputPos}
\difficulty{3}
Let a three-work-tape Turing-machine configuration have input head position $p$, work
tapes $T_0,T_1,T_2$, work-head positions $h_0,h_1,h_2$, and output string $O$. A control
action that moves the input head by $im$, makes no changes to the work tapes or their
head positions, emits the optional bit $o$, and sets the state to $q$ produces the
configuration with state $q$, input head at the clamped position obtained by moving $p$
by $im$, the same work tapes and work-head positions, and output $O$ followed by $o$ if
present.
\end{theorem}

\begin{theorem}[Work-tape printing walk]
\lean{BoolCircuit.CktSatReduction.Emitter.print_walk}
\label{BoolCircuit.CktSatReduction.Emitter.print_walk}
\leanok
\uses{BoolCircuit.CktSatReduction.Emitter.St, BoolCircuit.CktSatReduction.Emitter.apply_act, BoolCircuit.CktSatReduction.Emitter.emTM, BoolCircuit.CktSatReduction.Emitter.step_some, BoolCircuit.CktSatReduction.Emitter.tr, BoolCircuit.GateKind, Turing.Cfg, Turing.MultiTapeTM.runFrom, Turing.MultiTapeTM.runFrom_succ_eq_step, Turing.MultiTapeTM.runFrom_zero, Turing.UnaryTape.ones, Turing.UnaryTape.ones_lt, Turing.UnaryTape.ones_neg, Turing.UnaryTape.wrT, Turing.moveInputPos_zero}
\difficulty{5}
Let $w$ be a Boolean input, let $k$ be a gate kind, $c$ a gate code, $pc$ a program
counter, and $t$ one of the three work tapes. Suppose the input-head position is $p$,
the work tapes are given by $tp$, their head positions by $hd$, and tape $t$ contains
the unary encoding of $m$. For every $j\le m$ and every output list $out$, the emitter
machine started in its printing state for tape $t$, with that tape's head at $j-1$ and
the other heads as in $hd$, reaches after $j+1$ steps its printing state for the same
tape and program counter, with that tape's head at zero, the same input-head position
and tape contents, and output $out$ followed by $j$ true bits.
\end{theorem}

\begin{theorem}[Backward scan on a unary tape]
\lean{BoolCircuit.CktSatReduction.Emitter.back_walk}
\label{BoolCircuit.CktSatReduction.Emitter.back_walk}
\leanok
\uses{BoolCircuit.CktSatReduction.Emitter.St, BoolCircuit.CktSatReduction.Emitter.apply_act, BoolCircuit.CktSatReduction.Emitter.apply_ctl, BoolCircuit.CktSatReduction.Emitter.emTM, BoolCircuit.CktSatReduction.Emitter.step_some, BoolCircuit.CktSatReduction.Emitter.tr, BoolCircuit.GateKind, Turing.Cfg, Turing.MultiTapeTM.runFrom, Turing.MultiTapeTM.runFrom_succ_eq_step, Turing.MultiTapeTM.runFrom_zero, Turing.UnaryTape.ones, Turing.UnaryTape.ones_ge, Turing.UnaryTape.ones_lt, Turing.UnaryTape.wrT, Turing.moveInputPos_zero}
\difficulty{4}
Let $w$ be a finite Boolean input, let $p$ be an input-head position, and let $k$ be an
AND, OR, or NOT gate kind. Choose $c,t\in\{0,1,2\}$, a program counter
$pc\in\{0,\ldots,63\}$, three work tapes with head positions, a natural number $m$, and
an output word. Suppose tape $t$ contains $m$ consecutive ones starting at cell $0$. For
every natural number $d\le m$, the machine started in state $\mathrm{prB}(k,c,pc,t)$
with the head of tape $t$ at cell $m-d$ reaches state $\mathrm{em}(k,c,pc+1)$ after
$d+1$ steps, with the input-head position, all tape contents, all other head positions,
and the output word unchanged, and the head of tape $t$ at cell $m$.
\end{theorem}

\begin{theorem}[Clearing a unary tape]
\lean{BoolCircuit.CktSatReduction.Emitter.clear_walk}
\label{BoolCircuit.CktSatReduction.Emitter.clear_walk}
\leanok
\uses{BoolCircuit.CktSatReduction.Emitter.St, BoolCircuit.CktSatReduction.Emitter.apply_act, BoolCircuit.CktSatReduction.Emitter.clNext, BoolCircuit.CktSatReduction.Emitter.emTM, BoolCircuit.CktSatReduction.Emitter.step_some, BoolCircuit.CktSatReduction.Emitter.tr, Turing.Cfg, Turing.MultiTapeTM.runFrom, Turing.MultiTapeTM.runFrom_succ_eq_step, Turing.MultiTapeTM.runFrom_zero, Turing.UnaryTape.ones, Turing.UnaryTape.ones_lt, Turing.UnaryTape.ones_neg, Turing.UnaryTape.update_ones_pred, Turing.UnaryTape.wrT, Turing.moveInputPos_zero}
\difficulty{5}
Let $w$ be a finite Boolean input, let $p$ be an input-head position, and let
$T_0,T_1,T_2 : \mathbb{Z} \to \mathrm{Option}(\mathrm{Bool})$ and $h_0,h_1,h_2 \in
\mathbb{Z}$ be the contents and head positions of three work tapes. Choose an output
word $O$ and a tape index $t \in \{0,1,2\}$. For every $j \in \mathbb{N}$, if the
machine starts in state $\mathrm{clW}(t)$ with tape $t$ containing a $1$ at exactly the
cells $0,\ldots,j-1$ and its head at $j-1$, then after $j+1$ steps it is in state
$\mathrm{clS2}$ if $t=1$ and in state $\mathrm{inc}$ otherwise; tape $t$ is blank
everywhere with its head at $0$, while the input-head position, the other work tapes and
their head positions, and the output word remain unchanged.
\end{theorem}
